\documentclass[journal, 10pt, twocolumn]{IEEEtran}

% --- Essential Packages ---
\usepackage{xeCJK} % Required for Traditional Chinese support (Compile with XeLaTeX)
\usepackage{amsmath,amssymb,amsfonts}
\usepackage{algorithmic}
\usepackage{graphicx}
\usepackage{textcomp}
\usepackage{xcolor}
\usepackage{booktabs}
\usepackage{hyperref}
\usepackage{tikz}
\usetikzlibrary{shapes.geometric, arrows, positioning, calc, shadows}

% --- TikZ Styles ---
\tikzstyle{layer} = [rectangle, rounded corners, minimum width=3cm, minimum height=1cm, text centered, draw=black, fill=gray!5, font=\small]
\tikzstyle{core} = [rectangle, rounded corners, minimum width=3cm, minimum height=1cm, text centered, draw=green!60!black, fill=green!10, text=green!30!black, line width=1pt, font=\small]
\tikzstyle{arrow} = [thick,->,>=stealth]

\begin{document}

\title{Saccade：針對邊緣 AI 的高效率雙軌視覺感知系統}

\author{
    \IEEEauthorblockN{Saccade 專案團隊} \\
    \IEEEauthorblockA{\textit{邊緣運算與電腦視覺研究小組}}
}

\maketitle

% --- 摘要 (Abstract) ---
\begin{abstract}
邊緣運算設備（如僅具備 4GB VRAM 之設備）在執行連續視覺感知時面臨嚴重的資源瓶頸。本文提出 Saccade 系統，其核心論點聚焦於兩大面向：首先，針對\textbf{邊緣運算設備優化}，透過「零拷貝 (Zero-Copy)」硬體管線與原生 TensorRT 引擎，消除了資料搬運延遲，實現了 120 FPS 的高吞吐量；其次，整合了\textbf{視覺與語義搜尋的向量資料庫}，利用 SigLIP 2 模型與 ChromaDB 實現自然語言檢索。實驗數據顯示，系統能將單幀傳輸延遲壓縮至 0.00016 ms，並在 1.5GB VRAM 限制下穩定運作。
\end{abstract}

\begin{IEEEkeywords}
邊緣 AI, 零拷貝, 向量資料庫, 語義搜尋, Saccade.
\end{IEEEkeywords}

% --- 第一章：緒論 (Introduction) ---
\section{緒論}
\IEEEPARstart{隨}{著}邊緣 AI 的普及，如何在受限資源下實現「感知快」與「搜尋準」成為挑戰。Saccade 系統基於兩大核心論點進行開發：
\begin{itemize}
    \item \textbf{邊緣設備硬體管線優化：} 透過 Zero-Copy 與原生 C++ 遷移，解決 PCIe 頻寬與 GIL 瓶頸。
    \item \textbf{語義搜尋向量資料庫：} 建立具備長期語義記憶的視覺監控系統。
\end{itemize}

% --- 第二章：系統架構 (System Architecture) ---
\section{系統架構}
Saccade 採用六層式結構 (L1-L6)，結合了高性能感知與深度語義存儲（圖 \ref{fig:arch}）。

\begin{figure}[h]
\centering
\begin{tikzpicture}[node distance=1.5cm]
    \node (ingest) [core] {L0-L1: 邊緣硬體優化 (零拷貝)};
    \node (l2) [core, below=0.8cm of ingest] {L2: 語義提取 (SigLIP 2)};
    \node (l4) [layer, below=0.8cm of l2, fill=blue!5, draw=blue!50!black] {L4: 向量資料庫 (ChromaDB)};
    \node (l5) [layer, right=1.2cm of l4, fill=purple!5, draw=purple!50!black, minimum width=2.5cm] {L5: 語義搜尋介面};
    
    \draw [arrow] (ingest) -- (l2);
    \draw [arrow] (l2) -- (l4);
    \draw [arrow] (l5) -- (l4);
\end{tikzpicture}
\caption{Saccade 系統核心架構圖}
\label{fig:arch}
\end{figure}

% --- 第三章：核心技術實作 ---
\section{核心技術實作}
\subsection{Zero-Copy 與 5-Buffer GPU Pool}
透過 PyBind11 與 C++，影像解碼後直接作為 CUDA Tensor 駐留。單幀傳輸延遲僅 \textbf{0.00016 ms}，幾乎消除了傳統 numpy 轉換的開銷。

\subsection{語義過濾 (Semantic Drift)}
僅當物體發生「意義變化」時才寫入資料庫，過濾吞吐量達每秒 9,700+ 個物件，有效減少了儲存膨脹。

% --- 第四章：效能評估 (Performance Evaluation) ---
\section{實驗與效能評估}
我們在具備 \textbf{NVIDIA GeForce RTX 5070 Ti Laptop GPU (12GB)} 的設備上模擬 4GB 邊緣環境進行測試。

\subsection{細分延遲測試 (Latency Breakdown)}
下表展示了感知管線中各階段的處理時間（表 \ref{tab:latency}）。
\begin{table}[h]
\caption{各階段處理延遲}
\label{tab:latency}
\centering
\begin{tabular}{@{}lc@{}}
\toprule
處理階段 (Stage) & 平均耗時 (ms) \\ \midrule
Zero-Copy 資料傳輸 & \textbf{0.00016} \\
L1: YOLO 推理 (TensorRT) & 3.87 \\
L2: ROI Align (GPU 裁切) & 0.29 \\
L2: SigLIP 2 特徵提取 & 1.45 \\ \midrule
\textbf{端到端 (E2E) 總延遲} & \textbf{8.31} \\ \bottomrule
\end{tabular}
\end{table}

\subsection{吞吐量與資源佔用 (Throughput \& Resources)}
系統在維持高幀率的同時，成功將資源消耗壓制在預定義的臨界值內（表 \ref{tab:resources}）。
\begin{table}[h]
\caption{系統吞吐量與資源指標}
\label{tab:resources}
\centering
\begin{tabular}{@{}lcc@{}}
\toprule
指標項目 & 測試結果 & 單位 \\ \midrule
Fast Path 吞吐量 & \textbf{120.2} & FPS \\
語義過濾吞吐量 & 9,773 & objects/s \\
VRAM 佔用 (全負載) & 1.42 & GB \\
CPU 使用率 & < 15 & \% \\ \bottomrule
\end{tabular}
\end{table}

\subsection{儲存與檢索效率 (Storage \& Search)}
向量資料庫的效能支撐了智慧搜尋的即時性。
\begin{itemize}
    \item \textbf{資料庫寫入延遲：} 220 ms/record (透過 Redis Streams 非同步緩衝)。
    \item \textbf{語義搜尋延遲：} 206 ms (針對大量向量進行搜尋)。
    \item \textbf{存儲空間節省：} > 90\% (得益於語義過濾機制)。
\end{itemize}

% --- 第五章：結論 ---
\section{結論}
Saccade 驗證了透過極致硬體優化與語義搜尋架構，邊緣 AI 也能具備高速且具備「理解力」的視覺感知。這套數據證實了該方案在實際部署中的可行性與優越性。

\end{document}